% Generated draft: requires mathematical and source review.
\documentclass[11pt]{article}
\usepackage[margin=25mm]{geometry}
\usepackage{fontspec}
\setmainfont{Latin Modern Roman}
\usepackage{amsmath,amssymb,mathtools,unicode-math}
\setmathfont{Latin Modern Math}
\usepackage{enumitem,needspace,xurl,hyperref}
\setlength{\parindent}{0pt}
\setlength{\parskip}{5pt}
\setlength{\emergencystretch}{3em}
\newcommand{\N}{\mathbb{N}}
\newcommand{\Z}{\mathbb{Z}}
\newcommand{\Q}{\mathbb{Q}}
\newcommand{\R}{\mathbb{R}}
\newcommand{\C}{\mathbb{C}}
\newcommand{\F}{\mathbb{F}}
\newcommand{\A}{\mathbb{A}}
\newcommand{\PP}{\mathbb{P}}
\newcommand{\E}{\mathbb{E}}
\newcommand{\eps}{\varepsilon}
\newcommand{\dd}{\,\mathrm{d}}
\begin{document}
\section*{UnsolvedMath ID: 701152}
\textbf{Source problem code:} catalogue entry
\subsection{Definitions and mathematical statement}
Let \(G\) and \(X\) be finite groups. An action of \(G\) on \(X\) by group automorphisms is a homomorphism \(\rho:G\to\operatorname{Aut}(X)\). Two such actions \(\rho\) and \(\sigma\) are equivalent if there exists \(\alpha\in\operatorname{Aut}(X)\) such that
\[
\sigma(g)=\alpha\rho(g)\alpha^{-1}\qquad\text{for every }g\in G.
\]
Define
\[
r(G;X)=\left|\operatorname{Hom}(G,\operatorname{Aut}(X))/\operatorname{Aut}(X)\right|,
\qquad
r_G(n)=\max_{\substack{X\text{ a finite group}\\|X|=n}}r(G;X)
\quad(n\geq1),
\]
where the quotient denotes the orbits under conjugation, and the maximum may be taken over isomorphism classes of groups of order \(n\). For the prime factorization \(n=\prod_{i=1}^{t}p_i^{a_i}\), with distinct primes \(p_i\) and positive integers \(a_i\), set
\[
\lambda(n)=\sum_{i=1}^{t}a_i,\qquad\lambda(1)=0.
\]
Under the interpretation that the bound may depend on the fixed finite group \(G\), the question of P. M. Neumann is whether
\[
\forall\,G\text{ finite}\quad
\exists\,f_G:\Z_{\geq0}\to\Z_{\geq0}\quad
\forall\,n\in\Z_{\geq1},\qquad
r_G(n)\leq f_G(\lambda(n)).
\]
The supplied wording does not explicitly specify whether dependence on \(G\) is permitted. A bound uniform in \(G\) would instead require a single function \(f\) preceding the quantifier over finite groups \(G\); these interpretations are distinct, and the intended dependence requires review.
\subsection{Short English statement}
For a fixed finite group \(G\), count actions by automorphisms on each finite group \(X\), identifying actions conjugate by an automorphism of \(X\). Is the largest count among groups of order \(n\) bounded in terms of the number of prime factors of \(n\), counted with multiplicity? The formulation above permits dependence on \(G\); the source does not explicitly settle that qualification.
\subsection{Sources}
\begin{itemize}
\item Local UnsolvedMath registry, permanent ID \(701152\), exact source problem code \texttt{catalogue entry}, title \emph{Kourovka Notebook Problem 15.64}; supplied source record, upstream revision \texttt{d371ee917ac9647fbb06cd6fc529deec6f3dda82}.
\item P. M. Neumann, problem 15.64, \emph{The Kourovka Notebook (Unsolved Problems in Group Theory)}, 15th issue (2002). The supplied record also identifies the 21st issue edition (2026), \texttt{arXiv:1401.0300v48}. The collection description names E. I. Khukhro and V. D. Mazurov as editors.
\item The registry labels this problem \texttt{open}; its supplied update date is \texttt{2026-10-16T14:00:00Z}. This is supplied registry metadata, not an independently verified literature assessment. The cited notebook text has not been independently consulted for this draft.
\end{itemize}
\end{document}
